% !TeX program = lualatex
% !TeX encoding = utf8
% !TeX spellcheck = uk_UA
% !TeX root =../ClassicalElectrodynamics.tex

% chktex-file 3
% chktex-file 21
% chktex-file 25
% chktex-file 46

%% --------------------------------------------------------
\section{Узагальнені функції}\setcounter{equation}{0}\label{sec:GeneralizedFunctions}
%% --------------------------------------------------------

%Дельта-функція Дірака (або $\delta$-функція) є узагальненою функцією і була введена фізиком Полем Діраком для моделювання густини ідеалізованої
%точкової маси або точкового заряду.
%
%На <<фізичному рівні строгості>> можна визначити $\delta$-функцію формальним співвідношенням:
%\begin{equation}
%    \int\limits_{-\infty}^{+\infty}f(x)\delta(x-x_0)\,dx=f(x_0),
%\end{equation}
%У випадку інтегрування по скінченному об'єму $V$:
%\begin{equation}\label{eq:delta3D}
%    \int\limits_{V}f(\vect{r})\delta(\vect{r} - \vect{r}_0) dV=f(\vect{r}_0),
%\end{equation}
%де точка $\vect{r}_0$ знаходиться всередині об'єму $V$.

%$\delta$-функцію однієї дійсної змінної можна визначити як функцію, що задовольняє наступним умовам:
%\begin{equation}\label{eq:delta}
%\delta(x - x_0)=\left\{\begin{matrix}
%   +\infty, & x=x_0, \\
%   0, & x \neq x_0; \\
%\end{matrix}\right.
%\end{equation}
%\begin{equation}\label{eq:deltaProp0}
%    \int\limits_{-\infty}^{+\infty}\delta(x-x_0)dx=1.
%\end{equation}
%Тобто ця функція не дорівнює нулю тільки в точці $x = x_0$, де вона перетворюється в нескінченність таким чином, щоб її інтеграл в будь-якому околі точки $x = x_0$ дорівнює $1$.
%
%Аналітичне представлення $\delta$-функції:
%\begin{equation}\label{eq:deltaanalit}
%    \delta(x) = \frac{1}{2\pi}\int\limits_{-\infty}^{\infty} e^{ikx}dk.
%\end{equation}
%\bigskip\noindent%
%\textbf{Властивості дельта-функції}
%\bigskip
%
%\begin{enumerate}[label=\alph*)]
%\item Дельта-функція парна $\delta(-x) = \delta(x)$,
%\item $x\delta(x) = 0$,
%\item $\delta(ax) = \frac{1}{|a|}\delta(x)$,
%\item  $x\delta^\prime(x)=-\delta(x)$,
%\item $\delta(f(x))=\sum\limits_k\frac{\delta(x-x_k)}{|f'(x_k)|}$, де $x_k$~--- нулі функції $f(x)$,
%
%\end{enumerate}



В задачах електродинаміки часто виникає необхідність розглядати заряджені тіла, розміри яких дуже малі у порівнянні з іншими просторовими масштабами. З
цим пов'язана модель точкового (сферичного) заряду, густину якого подають у вигляді:

\begin{equation}
\rho(\vect{r}) = q \delta(\vect{r} - \vect{r}_0),
\end{equation}
%
де \( q \) --- величина заряду, \( \vect{r} \) --- його положення, \( \delta \) --- функція Дірака, що визначається умовою:

\begin{equation}
\int dV \, \delta(\vect{r} - \vect{r}_0) \chi(\vect{r}) = \chi(\vect{r}_0)
\label{eq:dirac_definition}
\end{equation}
%
для будь-якої неперервної функції \( \chi \).

Очевидно, що визначення~\eqref{eq:dirac_definition} \(\delta\)-функції не є математично коректним, якщо вважати, що в~\eqref{eq:dirac_definition} маємо
справу зі звичайним інтегруванням, наприклад, за Лебегом. Цьому співвідношенню можна надати математичний зміст, якщо розглядати \(\delta\)-функцію як
слабку границю деякої послідовності звичайних функцій. Якщо покласти (в одновимірному варіанті):

\begin{equation}
\delta_{\varepsilon}(x) = \frac{1}{\varepsilon \sqrt{\pi}} e^{-x^2 / \varepsilon^2},
\label{eq:delta_sequence}
\end{equation}

то границя:

\begin{equation}
\lim_{\varepsilon \to 0} \int dx \, \delta_{\varepsilon}(x) \chi(x) = \chi(0)
\label{eq:delta_limit}
\end{equation}
існує для будь-якої послідовності \( \varepsilon_n \to 0 \), \( n \to \infty \) та будь-якої функції \( \chi(x) \), що задовольняє досить широким умовам
(наприклад, якщо \( \chi(x) \) --- гладка та обмежена). На практиці зручніше розглядати подібні границі за жорсткіших обмежень на \( \chi(x) \).

Таке визначення цілком відповідає фізичним уявленням про точковий заряд, зосереджений в нескінченно малій області. Насправді для багатьох застосувань
істотно лише те, що розміри заряду значно менші за відстань до нього. Разом із тим треба пам'ятати, що ці міркування не проходять, коли треба розглядати
співвідношення, нелінійні за густиною заряду або за напруженостями полів; наприклад, у формулах для енергії.

%% --------------------------------------------------------
\subsection{Узагальнені та основні функції (\texorpdfstring{\( x \in \mathbb{R}^n \)}{x ∈ Rⁿ})}
%% --------------------------------------------------------

Узагальнену \(\delta\)-функцію можна визначити за допомогою іншої, відмінної від~\eqref{eq:delta_sequence}, послідовності, тобто означення
\eqref{eq:delta_sequence}, \eqref{eq:delta_limit} не є єдиним (див. приклади далі). Але усі подібні співвідношення визначають лінійний неперервний
функціонал, що зіставляє функції \( \chi(x) \) число \( \chi(0) \). Далі розглянемо лінійні неперервні функціонали на деякій множині функцій, які будемо
називати \textit{основними}.

Нехай множина основних функцій --- це множина \( D \) фінітних (тобто відмінних від нуля у деяких обмежених областях) нескінченно диференційовних
функцій. Збіжність в \( D \) визначають так: послідовність \( \{\phi_k\} \subset D \) збігається до \( \phi \in D \), якщо всі \( \phi_k \) можуть бути
відмінними від нуля лише у деякій спільній обмеженій області та усі похідні рівномірно збігаються до відповідних похідних від \( \phi \).

Узагальнена функція --- це лінійний неперервний функціонал, означений на \( D \). Множину таких узагальнених функцій позначають через \( D' \).

Існує багато способів регуляризації дельта-функції Дірака; зокрема, використовуються наступні наближення:

\begin{equation}
    \delta_{\varepsilon}(x) = \frac{1}{\pi} \frac{\sin(x/\varepsilon)}{x},
\label{eq:delta_sin}
\end{equation}

\begin{equation}
    \delta_{\varepsilon}(x) = \frac{1}{\pi} \frac{\varepsilon}{x^2 + \varepsilon^2},
\label{eq:delta_lorentz}
\end{equation}

\begin{equation}
    \delta_{\varepsilon}(x) = \frac{1}{\pi} \frac{\sin^2(x/\varepsilon)}{x^2},
\label{eq:delta_sin2}
\end{equation}
де \( x \in \mathbb{R} \).

Далі будемо позначати значення функціоналу, що відповідає узагальненій функції \( F \in D' \) на основній функції \( \chi \in D \), як \( (F, \chi) \).
Наприклад, співвідношення \( (\delta, \chi) = \chi(0) \) визначає \(\delta\)-функцію Дірака. Узагальнену функцію \( F(x) \) (\( x \in \mathbb{R}^n \))
називають \textit{регулярною}, якщо існує інтегровна функція \( f(x) \), така, що відповідний функціонал можна подати у вигляді звичайного інтеграла
Лебега:

\begin{equation}\label{eq:regular_functional}
    (F, \chi) = \int dx \, f(x) \chi(x).
\end{equation}

Якщо це неможливо, \( F \) називають \textit{сингулярною узагальненою функцією}. Прикладом такої функції є \(\delta\)-функція Дірака. Проте, у
фізичній літературі для сингулярних функцій також використовують запис~\eqref{eq:regular_functional}, маючи на увазі певний граничний перехід --- типу
\eqref{eq:delta_limit} або інший. Можна показати, що будь-яку узагальнену функцію можна подати як слабку границю послідовності основних функцій з~\( D
\).

Далі ми також будемо застосовувати формальний запис~\eqref{eq:regular_functional}, пам'ятаючи про вказані застереження.

%% --------------------------------------------------------
\subsection{Диференціювання і фундаментальні розв'язки}
%% --------------------------------------------------------

Визначимо похідну \( \partial_i F(x) \) від узагальненої функції \( F \) співвідношенням:

\begin{equation}\label{eq:derivative_definition}
    (\partial_i F, \chi) = - (F, \partial_i \chi), \quad x \in \mathbb{R}^n.
\end{equation}

Для регулярної диференційовної функції \( F \) це співвідношення є очевидним наслідком інтегрування за частинами:

\begin{equation*}
    \int dx \, \partial_i F(x) \chi(x) = - \int dx \, F(x) \partial_i \chi(x),
\end{equation*}
де межові члени зникають, оскільки \( \chi(x) \) є фінітною.

Нехай \( \hat{L} \) --- диференціальний оператор скінченого порядку зі сталими коефіцієнтами, \( x \in \mathbb{R}^n \).

\emph{Фундаментальним розв'язком оператора \( \hat{L} \)} називають узагальнену функцію \( G \), таку, що:

\begin{equation}\label{eq:fundamental_solution}
    \hat{L} G = \delta(x).
\end{equation}

%% --------------------------------------------------------
\subsubsection{Фундаментальний розв'язок оператора Лапласа}
%% --------------------------------------------------------

Нехай  \( x = \vect{r} = (x_1, x_2, x_3) \in \mathbb{R}^3 \), покажемо, що:

\begin{equation}\label{eq:laplace_fundamental}
    \Laplasian \left( \frac{1}{r} \right) = -4\pi \delta(\vect{r}),
\end{equation}

де в правій частині фігурує тривимірна \(\delta\)-функція,

\begin{equation*}
r = |\vect{r}| = \sqrt{x_1^2 + x_2^2 + x_3^2}.
\end{equation*}

За означенням похідної~\eqref{eq:derivative_definition}, застосованим двічі (межові члени зникають через фінітність \(\chi\)):

\begin{equation*}
\int \Laplasian \left( \frac{1}{r} \right) \chi(\vect{r}) dV = \int \frac{1}{r} \Laplasian \chi dV.
\end{equation*}

Права частина означена як звичайний невласний інтеграл, який можна обчислювати у сферичних координатах:

\begin{multline*}
\int \frac{\Laplasian \chi}{r}  r^2 \sin \theta \, d\theta \, d\varphi \, dr = \\
=
\int \frac{1}{r} \left[ \frac{\partial}{\partial r} \left( r^2
\frac{\partial \chi}{\partial r} \right) + \frac{1}{\sin \theta} \frac{\partial}{\partial \theta} \left( \sin \theta \frac{\partial \chi}{\partial
\theta} \right) + \frac{1}{\sin^2 \theta} \frac{\partial^2 \chi}{\partial \varphi^2} \right] \sin \theta \, d\theta \, d\varphi =
\\=
\int \frac{1}{r}
\frac{\partial}{\partial r} \left( r^2 \frac{\partial \chi}{\partial r} \right) dr d\Omega, \quad d\Omega = \sin \theta \, d\theta \, d\varphi.
\end{multline*}

Тут враховано, що:

\begin{equation*}
\int\limits_0^\pi d\theta \, \frac{\partial}{\partial \theta} \left( \sin \theta \frac{\partial \chi}{\partial \theta} \right) = \left( \sin \theta
\frac{\partial \chi}{\partial \theta} \right)_0^\pi = 0, \quad \int\limits_0^{2\pi} d\varphi \, \frac{\partial^2 \chi}{\partial \varphi^2} =
\frac{\partial
\chi}{\partial \varphi} \bigg|_0^{2\pi} = 0.
\end{equation*}

Залишається інтеграл, що береться по \( r \) при фіксованих \( \varphi \), \( \theta \):

\begin{multline*}
\int\limits_0^\infty dr \int d\Omega \, \frac{1}{r} \frac{\partial}{\partial r} \left( r^2 \frac{\partial \chi}{\partial r} \right) =
\int d\Omega
\int\limits_0^\infty dr \, \frac{\partial}{\partial r} \left( r \frac{\partial \chi}{\partial r} + \chi \right) = \left.\left( r \frac{\partial
\chi}{\partial r} +
\chi
\right)\right|_0^\infty = \\ =- \int d\Omega \, \chi(0) = -4\pi \chi(0).
\end{multline*}

Отриманий вираз доводить співвідношення~\eqref{eq:laplace_fundamental}. Відповідно до загального означення~\eqref{eq:fundamental_solution}, фундаментальним розв'язком оператора Лапласа є функція

\begin{equation}\label{eq:laplace_G}
    G(\vect{r}) = -\frac{1}{4\pi r}, \qquad \Laplasian G = \delta(\vect{r}).
\end{equation}

%% --------------------------------------------------------
\subsubsection{Фундаментальний розв'язок оператора Даламбера (хвильового оператора)}
%% --------------------------------------------------------

Нехай \( c = 1 \), \( \Dalambertian = \frac{\partial^2}{\partial t^2} - \Laplasian \),
\( x = (t, x_1, x_2, x_3) \in \mathbb{R}^4 \).

Визначимо функцію \( G(x) \) співвідношенням:

\begin{equation}
    G(x) = \frac{1}{4\pi r} \delta(t - r),
    \label{eq:dalambert_fundamental}
\end{equation}
%
де \( r = |\vect{r}| = \sqrt{x_1^2 + x_2^2 + x_3^2} \).

Покажемо, що \( \Dalambertian G = \delta(x) \), тобто
\( (\Dalambertian G,\, \chi) = \chi(0,\vect{0}) \) для будь-якої \( \chi \in D \).

За означенням похідної узагальненої функції~\eqref{eq:derivative_definition}:
\begin{equation*}
    (\Dalambertian G,\, \chi) = (G,\, \Dalambertian \chi)
    = \int dt\, dV\; G(t,\vect{r})\, \Dalambertian \chi(t,\vect{r}).
\end{equation*}

Підставляємо~\eqref{eq:dalambert_fundamental} і беремо інтеграл по \( t \)
за допомогою \( \delta(t-r) \):
\begin{equation*}
    (G, \Dalambertian \chi)
    = \int \frac{dV}{4\pi r}\, \Dalambertian \chi \big|_{t=r}
    = \frac{1}{4\pi} \int_0^\infty dr\, r \int d\Omega\;
      \Dalambertian \chi(r, r,\theta,\varphi),
\end{equation*}
де \( d\Omega = \sin\theta\,d\theta\,d\varphi \), а першим аргументом \( \chi \) є час \( t = r \).

Перейдемо до сферичних координат. Оператор Даламбера розкладається як
\begin{equation*}
    \Dalambertian \chi = \chi_{tt} - \frac{1}{r^2}\partial_r(r^2 \partial_r \chi)
    - \frac{1}{r^2}\hat{L}^2_\Omega \chi,
\end{equation*}
де \( \hat{L}^2_\Omega \) --- кутова частина Лапласіана. Інтегрування кутових членів по замкненій сфері дає нуль з тих самих міркувань, що і в доведенні
для оператора Лапласа. Залишається:
\begin{equation*}
    (G, \Dalambertian \chi)
    = \frac{1}{4\pi}\int d\Omega \int_0^\infty r\,
      \left(\chi_{tt} - \frac{1}{r^2}\partial_r(r^2\partial_r \chi)\right)_{t=r} dr.
\end{equation*}

Введемо позначення \( \psi(r) \equiv \chi(r, r,\theta,\varphi) \) (функція по діагоналі \( t = r \)) та запишемо похідні вздовж цієї діагоналі:
\begin{equation*}
    \frac{d}{dr}\chi(t,r)\big|_{t=r} = \chi_t + \chi_r, \qquad
    \frac{d^2}{dr^2}\chi(t,r)\big|_{t=r} = \chi_{tt} + 2\chi_{tr} + \chi_{rr}.
\end{equation*}

Скористаємося тотожністю (яку можна перевірити прямим диференціюванням):
\begin{equation}
    r\bigl(\chi_{tt} - \chi_{rr}\bigr)\big|_{t=r} - 2\chi_r\big|_{t=r}
    = \frac{d}{dr}\Bigl\{ r\bigl(\chi_t - \chi_r\bigr)\big|_{t=r} - \chi\big|_{t=r} \Bigr\}.
    \label{eq:dalambert_identity}
\end{equation}

Водночас:
\begin{equation*}
    \frac{1}{r^2}\partial_r(r^2\partial_r \chi)\big|_{t=r}
    = \chi_{rr} + \frac{2}{r}\chi_r,
\end{equation*}
тому:
\begin{equation*}
    r\left(\chi_{tt} - \frac{1}{r^2}\partial_r(r^2\partial_r\chi)\right)_{t=r}
    = r(\chi_{tt} - \chi_{rr})\big|_{t=r} - 2\chi_r\big|_{t=r}
    = \frac{d}{dr}\Bigl\{ r(\chi_t - \chi_r)\big|_{t=r} - \chi\big|_{t=r} \Bigr\}.
\end{equation*}

Отже, підінтегральний вираз є повною похідною. Інтегруємо по \( r \) від \( 0 \) до \( +\infty \):
\begin{multline*}
    \int_0^\infty \frac{d}{dr}\Bigl\{ r(\chi_t - \chi_r)\big|_{t=r} - \chi\big|_{t=r} \Bigr\} dr \\
    = \Bigl[ r(\chi_t - \chi_r) - \chi \Bigr]_{t=r}\Bigg|_0^\infty
    = 0 - \bigl(0 - \chi(0,\vect{0})\bigr) = \chi(0,\vect{0}).
\end{multline*}
Тут при \( r \to \infty \) вираз звертається до нуля через фінітність \( \chi \), а при \( r = 0 \) перший доданок \( r(\chi_t - \chi_r) \to 0 \).

Інтегруючи по кутах \( \int d\Omega = 4\pi \), отримуємо:
\begin{equation*}
    (G, \Dalambertian \chi) = \frac{1}{4\pi} \cdot 4\pi \cdot \chi(0,\vect{0}) = \chi(0,\vect{0}),
\end{equation*}
що і доводить \( \Dalambertian G = \delta(x) \).

%% --------------------------------------------------------
\subsubsection{Фундаментальний розв'язок оператора Гельмгольца}
%% --------------------------------------------------------

Покажемо, що фундаментальним розв'язком оператора Гельмгольца
\( \hat{L} = \Laplasian + k^2 \) (\( k \in \mathbb{R} \), \( k \neq 0 \)) є функція:

\begin{equation}
    G(\vect{r}) = -\frac{e^{ikr}}{4\pi r}, \qquad r = |\vect{r}|,
    \label{eq:helmholtz_fundamental}
\end{equation}

тобто \( (\Laplasian + k^2) G = \delta(\vect{r}) \).

За означенням похідної узагальненої функції~\eqref{eq:derivative_definition} та лінійністю:
\begin{equation*}
    \bigl((\Laplasian + k^2)G,\, \chi\bigr)
    = \bigl(G,\, (\Laplasian + k^2)\chi\bigr)
    = \int dV\; G(\vect{r})\,(\Laplasian + k^2)\chi(\vect{r}).
\end{equation*}

Підставляємо~\eqref{eq:helmholtz_fundamental}:
\begin{equation*}
    (G,\, (\Laplasian + k^2)\chi)
    = -\frac{1}{4\pi}\int \frac{e^{ikr}}{r}\,(\Laplasian + k^2)\chi\, dV.
\end{equation*}

Інтеграл розбиваємо на дві частини:
\begin{equation}
    I = I_1 + I_2, \qquad
    I_1 = -\frac{1}{4\pi}\int \frac{e^{ikr}}{r}\,\Laplasian\chi\, dV, \qquad
    I_2 = -\frac{k^2}{4\pi}\int \frac{e^{ikr}}{r}\,\chi\, dV.
    \label{eq:helm_split}
\end{equation}

Застосуємо формулу Гріна до \( I_1 \):
\begin{equation*}
    \int_V \bigl( u\,\Laplasian v - v\,\Laplasian u \bigr) dV
    = \oint_{\partial V} \bigl( u\,\nabla v - v\,\nabla u \bigr)\cdot d\vect{S}.
\end{equation*}

Виберемо кулю \( B_\varepsilon = \{r \leq \varepsilon\} \) і доповнення до неї \( V_\varepsilon = \mathbb{R}^3 \setminus B_\varepsilon \).
На \( V_\varepsilon \) функція \( u = e^{ikr}/r \) є гладкою і задовольняє рівнянню Гельмгольца:
\begin{equation*}
    \Laplasian \frac{e^{ikr}}{r} = -k^2 \frac{e^{ikr}}{r}, \quad r > 0.
\end{equation*}

Застосовуємо формулу Гріна на \( V_\varepsilon \):
\begin{equation*}
    \int_{V_\varepsilon} \frac{e^{ikr}}{r}\,\Laplasian\chi\, dV
    = \int_{V_\varepsilon} \chi\, \Laplasian\frac{e^{ikr}}{r}\, dV
    + \oint_{r=\varepsilon} \left( \frac{e^{ikr}}{r}\,\partial_r\chi
    - \chi\,\partial_r\frac{e^{ikr}}{r} \right) dS.
\end{equation*}

Перший інтеграл в правій частині дає \( -k^2\int_{V_\varepsilon} (e^{ikr}/r)\,\chi\, dV \), тому:
\begin{equation*}
    I_1 = -\frac{1}{4\pi}\left[
        -k^2\int_{V_\varepsilon} \frac{e^{ikr}}{r}\chi\, dV
        + \oint_{r=\varepsilon} \left( \frac{e^{ikr}}{r}\partial_r\chi
        - \chi\,\partial_r\frac{e^{ikr}}{r} \right) dS
    \right] + O(\varepsilon).
\end{equation*}

Підставляємо в~\eqref{eq:helm_split}: члени з \( k^2 \) скорочуються з \( I_2 \) (у границі \( \varepsilon \to 0 \)), і залишається лише поверхневий
інтеграл по сфері \( r = \varepsilon \):
\begin{equation*}
    I = \frac{1}{4\pi}\oint_{r=\varepsilon}\left(
        \frac{e^{ikr}}{r}\partial_r\chi - \chi\,\partial_r\frac{e^{ikr}}{r}
    \right) dS.
\end{equation*}

На сфері \( r = \varepsilon \) маємо \( dS = \varepsilon^2\,d\Omega \),
а при \( \varepsilon \to 0 \): \( e^{ik\varepsilon} \to 1 \), \( \partial_r(e^{ikr}/r)\big|_{r=\varepsilon} = (ik\varepsilon - 1)/\varepsilon^2 + O(1) \).
Тому:
\begin{equation*}
    I = \frac{1}{4\pi}\oint d\Omega\left(
        \varepsilon\,\partial_r\chi\big|_{r=\varepsilon}
        - \chi\big|_{r=\varepsilon}\,(ik\varepsilon - 1)
    \right) + O(\varepsilon).
\end{equation*}

При \( \varepsilon \to 0 \) перший доданок прямує до нуля, а другий дає:
\begin{equation*}
    I \xrightarrow{\varepsilon\to 0}
    \frac{1}{4\pi}\int d\Omega\;\chi(0) = \frac{4\pi}{4\pi}\,\chi(0) = \chi(0).
\end{equation*}

Таким чином, \( \bigl((\Laplasian + k^2)G,\chi\bigr) = \chi(0) \), що доводить~\eqref{eq:helmholtz_fundamental}.

%% --------------------------------------------------------
\subsubsection{Згортка}
%% ---------------------------------------------------

Згортка \( \zeta = F * \chi \) узагальненої функції \( F \) з основною \( \chi \) є за визначенням:

\begin{equation}
\zeta(y) = \int dx \, F(x) \chi(y - x).
\label{eq:convolution}
\end{equation}

% Щоб отримати розв'язок рівняння $\hat{L}\phi=f$ для довільної (досить регулярної) правої частини $f$, потрібно \enquote{просумувати} внески від точкових джерел, розподілених із густиною $f$ --- саме цю операцію і виконує згортка \( \zeta = F * \chi \) узагальненої функції \( F \) з основною \( \chi \) є за визначенням:


Знайдений фундаментальний розв'язок $G$ сам по собі відповідає точковому джерелу $\delta(x)$. За допомогою згортки та фундаментального розв'язку можна будувати розв'язки рівнянь виду:

\begin{equation}
\hat{L} \phi = f, \quad f \in D.
\end{equation}

Дійсно, якщо покласти \( \phi = G * f \), де \( \hat{L} G = \delta(x) \), маємо:

\begin{equation*}
\hat{L}\phi(y) = \int dx \, G(x) \hat{L}_y f(y - x) = \int dx \, \hat{L}_x G(x) f(y - x) = (\delta(x), f(y - x)) = f(y).
\end{equation*}
Тут введено індекси, щоб підкреслити, що \( \hat{L}_x \) діє на змінну \( x \), \( \hat{L}_y \) --- на \( y \). Перехід \( \hat{L}_y f(y - x) = \hat{L}_x f(y - x) \) використовує те, що \( \hat{L} \) має сталі коефіцієнти: похідна по \( y_i \) від \( f(y-x) \) збігається з похідною по \( x_i \). Цей результат за певних умов можна розширити на випадок \( f \in D' \).

Звідси маємо такі розв'язки рівнянь:

\begin{equation}
\Laplasian \phi = -4\pi \rho \implies \phi = \int dV' \, \frac{\rho(\vect{r}')}{|\vect{r} - \vect{r}'|}, \quad \phi = \phi(\vect{r}),
\end{equation}

\begin{equation}
\Dalambertian \phi = 4\pi \rho \implies \phi = \int \frac{\rho(t - |\vect{r} - \vect{r}'|, \vect{r}')}{|\vect{r} - \vect{r}'|} dV' ,
\quad
\phi = \phi(t, \vect{r}),
\end{equation}

\begin{equation}
\Laplasian \phi + k^2 \phi = -4\pi \rho \implies \phi(\vect{r}) = \int \frac{e^{ik|\vect{r} - \vect{r}'|}}{|\vect{r} - \vect{r}'|}
\rho(\vect{r}')\,
dV'.
\end{equation}